\begin{ex}
Một hình tròn có bán kính là 5 cm. Đường kính của hình tròn đó là bao nhiêu xăng-ti-mét?
\shortans{10}
\loigiai{
Đường kính gấp đôi bán kính: $5 \times 2 = 10$ cm.
}
\end{ex}

\begin{ex}
Một hình tròn có bán kính là 8 cm. Đường kính của hình tròn đó là bao nhiêu xăng-ti-mét?
\shortans{16}
\loigiai{
$8 \times 2 = 16$ cm.
}
\end{ex}

\begin{ex}
Một hình tròn có đường kính là 14 cm. Bán kính của hình tròn đó là bao nhiêu xăng-ti-mét?
\shortans{7}
\loigiai{
$14 : 2 = 7$ cm.
}
\end{ex}

\begin{ex}
Một hình tròn có đường kính là 18 cm. Bán kính của hình tròn đó là bao nhiêu xăng-ti-mét?
\shortans{9}
\loigiai{
$18 : 2 = 9$ cm.
}
\end{ex}

\begin{ex}
Một hình tròn có bán kính là 12 cm. Đường kính của hình tròn đó là bao nhiêu xăng-ti-mét?
\shortans{24}
\loigiai{
$12 \times 2 = 24$ cm.
}
\end{ex}

\begin{ex}
Một hình tròn có đường kính là 26 cm. Bán kính của hình tròn đó là bao nhiêu xăng-ti-mét?
\shortans{13}
\loigiai{
$26 : 2 = 13$ cm.
}
\end{ex}

\begin{ex}
Trong một hình tròn, đường kính gấp mấy lần bán kính?
\shortans{2}
\loigiai{
Đường kính dài gấp 2 lần bán kính.
}
\end{ex}

\begin{ex}
Trong một hình tròn, tâm có phải là trung điểm của mọi đường kính không? (Nếu đúng ghi 1, sai ghi 0).
\shortans{1}
\loigiai{
Tâm hình tròn là trung điểm của mọi đường kính.
}
\end{ex}

\begin{ex}
Một compa có độ mở hai đầu bằng 6 cm để vẽ một đường tròn. Đường kính của hình tròn vẽ được là bao nhiêu xăng-ti-mét?
\shortans{12}
\loigiai{
Độ mở compa chính là bán kính: $R = 6$ cm $\Rightarrow$ Đường kính = $6 \times 2 = 12$ cm.
}
\end{ex}

\begin{ex}
Một mặt đồng hồ hình tròn có bán kính 10 cm. Đường kính của mặt đồng hồ đó là bao nhiêu xăng-ti-mét?
\shortans{20}
\loigiai{
$10 \times 2 = 20$ cm.
}
\end{ex}

\begin{ex}
Bánh xe đạp có đường kính là 60 cm. Bán kính của bánh xe đạp đó là bao nhiêu xăng-ti-mét?
\shortans{30}
\loigiai{
$60 : 2 = 30$ cm.
}
\end{ex}

\begin{ex}
Một chiếc đĩa tròn có đường kính 22 cm. Bán kính của chiếc đĩa đó là bao nhiêu xăng-ti-mét?
\shortans{11}
\loigiai{
$22 : 2 = 11$ cm.
}
\end{ex}

\begin{ex}
Hình tròn tâm O có bán kính OA = 7 cm. Đoạn thẳng OB cũng là một bán kính. Đoạn OB dài bao nhiêu xăng-ti-mét?
\shortans{7}
\loigiai{
Tất cả các bán kính trong hình tròn đều bằng nhau: $OB = OA = 7$ cm.
}
\end{ex}

\begin{ex}
Hình tròn tâm O có đường kính AB = 16 cm. Bán kính OM dài bao nhiêu xăng-ti-mét?
\shortans{8}
\loigiai{
$16 : 2 = 8$ cm.
}
\end{ex}

\begin{ex}
Nếu bán kính hình tròn tăng thêm 3 cm thì đường kính của hình tròn đó tăng thêm bao nhiêu xăng-ti-mét?
\shortans{6}
\loigiai{
$3 \times 2 = 6$ cm.
}
\end{ex}

\begin{ex}
Một hình tròn có đường kính là 40 cm. Bán kính của nó là bao nhiêu đề-xi-mét?
\shortans{2}
\loigiai{
$40 : 2 = 20\text{ cm} = 2\text{ dm}$.
}
\end{ex}

\begin{ex}
Một hình tròn có bán kính là 15 dm. Đường kính của hình tròn đó là bao nhiêu đề-xi-mét?
\shortans{30}
\loigiai{
$15 \times 2 = 30$ dm.
}
\end{ex}

\begin{ex}
Đường kính của hình tròn bằng 50 cm. Bán kính của hình tròn đó là bao nhiêu xăng-ti-mét?
\shortans{25}
\loigiai{
$50 : 2 = 25$ cm.
}
\end{ex}

\begin{ex}
Hình tròn có bao nhiêu tâm?
\shortans{1}
\loigiai{
Hình tròn chỉ có duy nhất 1 tâm.
}
\end{ex}

\begin{ex}
Hình tròn tâm O có bán kính 4 cm. Độ dài đường kính MN của hình tròn đó là bao nhiêu xăng-ti-mét?
\shortans{8}
\loigiai{
$4 \times 2 = 8$ cm.
}
\end{ex}

\begin{ex}
Biết bán kính một hình tròn bằng chu vi một hình vuông có cạnh 2 cm. Đường kính hình tròn đó là bao nhiêu xăng-ti-mét?
\shortans{16}
\loigiai{
Chu vi hình vuông = $2 \times 4 = 8$ cm. Bán kính = 8 cm $\Rightarrow$ đường kính = $8 \times 2 = 16$ cm.
}
\end{ex}

\begin{ex}
Độ mở hai đầu compa là 9 cm. Khi quay compa vẽ được hình tròn có đường kính bao nhiêu xăng-ti-mét?
\shortans{18}
\loigiai{
$9 \times 2 = 18$ cm.
}
\end{ex}

\begin{ex}
Đoạn thẳng nối tâm với một điểm trên hình tròn gọi là gì? (Bán kính hay đường kính? Nếu là bán kính ghi 1, đường kính ghi 2).
\shortans{1}
\loigiai{
Đoạn thẳng nối tâm với một điểm trên hình tròn là bán kính.
}
\end{ex}

\begin{ex}
Đoạn thẳng đi qua tâm và cắt hai điểm trên đường tròn gọi là gì? (Bán kính ghi 1, đường kính ghi 2).
\shortans{2}
\loigiai{
Đó là đường kính.
}
\end{ex}

\begin{ex}
Bán kính hình tròn là 35 cm. Đường kính hình tròn là bao nhiêu xăng-ti-mét?
\shortans{70}
\loigiai{
$35 \times 2 = 70$ cm.
}
\end{ex}
